\section{Механизм перехвата CUDA Runtime}

\subsection{Архитектура перехвата через \texttt{LD\_PRELOAD}}

\subsubsection{Схема взаимодействия компонентов}

Центральная идея механизма перехвата состоит в том, что CUDA-приложение
вызывает функции из \texttt{libcudart.so} через PLT, а динамический
загрузчик \texttt{ld.so} разрешает эти символы в реализации из
\texttt{libemuruntime.so}, заблаговременно помещённой в переменную
\texttt{LD\_PRELOAD}.

Полная схема взаимодействия компонентов:

\begin{figure}[h]
  \centering
  \includegraphics[width=16cm]{ldpreload_arch.png}
  \caption{Схема перехвата CUDA Runtime через \texttt{LD\_PRELOAD}.}
  \label{fig:ldpreload_arch}
\end{figure}

Библиотека \texttt{libemuruntime.so} реализует полный публичный ABI
\texttt{libcudart.so}, необходимый для работы типичного CUDA-приложения.
Каждая функция-заглушка либо делегирует работу компонентам эмулятора
(парсеру, исполнителю, профилировщику), либо возвращает \texttt{cudaSuccess}
для функций, не требующих реальной реализации в пределах эмулятора
(например, управление событиями, профили потоков).

\subsubsection{Реализованные точки перехвата CUDA Runtime API}

В таблице~\ref{tab:intercept_points} перечислены перехватываемые функции
CUDA Runtime API с указанием ответственного компонента.

\begin{table}[h]
  \caption{Реализованные точки перехвата CUDA Runtime API}
  \label{tab:intercept_points}
  \centering
  \begin{tabularx}{\linewidth}{lXl}
    \toprule
    \textbf{Функция} & \textbf{Назначение} & \textbf{Компонент} \\
    \midrule
    \texttt{\_\_cudaRegisterFatBinary}  & Регистрация модуля; вызов \texttt{cuobjdump} для извлечения PTX & Часть~1 \\
    \texttt{\_\_cudaRegisterFunction}   & Отображение дескриптор $\to$ имя ядра & Часть~1 \\
    \texttt{cudaLaunchKernel}           & Запуск ядра через диспетчер & Часть~2 \\
    \texttt{cudaMalloc}                 & Выделение памяти устройства & Часть~2 \\
    \texttt{cudaFree}                   & Освобождение памяти устройства & Часть~2 \\
    \texttt{cudaMemcpy}                 & Копирование памяти хост$\leftrightarrow$устройство & Часть~2 \\
    \texttt{cudaMemset}                 & Инициализация памяти & Часть~2 \\
    \texttt{cudaDeviceSynchronize}      & Синхронизация (no-op в эмуляторе) & Часть~2 \\
    \bottomrule
  \end{tabularx}
\end{table}

\subsubsection{Извлечение PTX-текста через \texttt{cuobjdump}}

При перехвате \texttt{\_\_cudaRegisterFatBinary(handle)} эмулятор
не разбирает fatbinary напрямую. Вместо этого он определяет путь
к исполняемому файлу приложения из переменной окружения
\texttt{CUEMU\_TARGET\_EXEC} и передаёт его утилите
\texttt{cuobjdump}:

\begin{lstlisting}[language=sh, caption={Извлечение PTX через cuobjdump}]
cuobjdump --dump-ptx <binary> > ptx_text.txt
\end{lstlisting}

\texttt{cuobjdump} извлекает PTX-текст из секции \texttt{.nv\_fatbin}
ELF-файла и выводит его в стандартный вывод. Эмулятор читает этот
вывод и передаёт PTX-текст в \texttt{PtxModule::Make()},
которая строит таблицу инструкций. Полученный \texttt{Module}
сохраняется в \texttt{GlobalContext} под ключом-дескриптором.

Алгоритм извлечения PTX:
\begin{enumerate}
  \item Прочитать путь к бинарному файлу из \texttt{CUEMU\_TARGET\_EXEC}.
  \item Выполнить \texttt{cuobjdump --dump-ptx <binary>},
        захватить stdout через \texttt{popen()}.
  \item Передать полученный текст в \texttt{PtxModule::Make()},
        получить объект \texttt{Module} с таблицей инструкций.
  \item Сохранить \texttt{Module} в \texttt{GlobalContext}
        под ключом-дескриптором.
\end{enumerate}

\subsubsection{Регистрация функций: отображение дескриптор~$\to$~имя ядра}

CUDA Runtime ABI определяет, что при вызове \texttt{cudaLaunchKernel}
ядро идентифицируется \textit{указателем на функцию-заглушку}
в сегменте кода, а не строковым именем. Это обусловлено тем, что
в C++ ядра компилируются в обычные функции, и CUDA не гарантирует
сохранение символьных имён после компоновки.

При перехвате \texttt{\_\_cudaRegisterFunction} эмулятор создаёт
отображение:
\begin{equation}
  \mathit{functions\_}[\mathtt{ptr}] = \mathit{name},
\end{equation}
где \texttt{ptr}~--- числовое представление указателя на заглушку,
\texttt{name}~--- строка вида \texttt{"\_Z11vadd\_kernelPfS\_S\_i"}.

При последующем вызове \texttt{cudaLaunchKernel(ptr, ...)} эмулятор
извлекает строковое имя из отображения по ключу \texttt{ptr}
и напрямую ищет функцию в \texttt{PtxModule::GetEntryFunction(name)}
без демангла~--- PTX-имя уже является декорированным символом,
которым \texttt{cuobjdump} подписывает каждую точку входа.

\subsubsection{Запуск \texttt{cuemu}: утилита командной строки и переменные окружения}

Утилита командной строки \texttt{cuemu} является точкой входа в программу: она подготавливает
окружение и повторно запускает целевой бинарный файл:

\begin{lstlisting}[language=sh, caption={Логика запуска cuemu}]
export LD_PRELOAD=$CUEMU_LIB_DIR/libemuruntime.so
export LD_LIBRARY_PATH=$CUEMU_LIB_DIR:$LD_LIBRARY_PATH
export CUEMU_TARGET_EXEC="$1"
export CUEMU_COLLECT_PROFILING="${CUEMU_COLLECT_PROFILING:-0}"
export CUEMU_PROFILING_OUTPUT="${CUEMU_PROFILING_OUTPUT:-profiling.txt}"
exec "$@"
\end{lstlisting}

Поддерживаемые переменные окружения перечислены в таблице~\ref{tab:env_vars}.

\begin{table}[h]
  \caption{Переменные окружения эмулятора}
  \label{tab:env_vars}
  \centering
  \begin{tabularx}{\linewidth}{lXl}
    \toprule
    \textbf{Переменная} & \textbf{Назначение} & \textbf{По умолчанию} \\
    \midrule
    \texttt{CUEMU\_COLLECT\_PROFILING} & Включить сбор метрик & \texttt{0} \\
    \texttt{CUEMU\_PROFILING\_OUTPUT}  & Путь к файлу вывода   & \texttt{profiling.txt} \\
    \texttt{CUEMU\_GPU\_CONFIG}        & Путь к TOML-конфигу GPU & встроенный \\
    \texttt{CUEMU\_TARGET\_EXEC}       & Имя целевого бинарного файла & обязателен \\
    \bottomrule
  \end{tabularx}
\end{table}

\newpage
